\subsection{表示定理的例子与证明核对}

\begin{example}[有限序的显式表示]
设 $X=\{a,b,c,d,e\}$ 且 $a\succ b\sim c\succ d\succ e$。把无差异类分为
$[a]=\{a\}$、$[b]=\{b,c\}$、$[d]=\{d\}$、$[e]=\{e\}$，依次赋值 $3,2,1,0$，则
$u(a)=3$、$u(b)=u(c)=2$、$u(d)=1$、$u(e)=0$。验证任意一对方案即可看到：严格顺序对应严格数值，不同类之间没有重复值，同一类内部数值相同。
\end{example}

\begin{proposition}[单调变换保持表示]
若 $u$ 表示 $\succeq$，$f$ 严格递增，则 $f\circ u$ 也表示 $\succeq$。
\end{proposition}
\begin{proof}
严格递增意味着对任意 $a,b$，$a\ge b$ 当且仅当 $f(a)\ge f(b)$。因此
$x\succeq y\Leftrightarrow u(x)\ge u(y)\Leftrightarrow f(u(x))\ge f(u(y))$，恰为表示定义。
\end{proof}

\begin{example}[同一偏好的三种写法]
对 $u(x_1,x_2)=x_1+x_2$，函数 $v=2u+1$ 与 $w=\ln(1+u)$ 在定义域上都严格递增变换，因此三者产生相同的无差异集合与最优选择。不能据此比较 $u(x)-u(y)$ 与 $v(x)-v(y)$ 的大小。
\end{example}

\subsection{可数构造的细节}
对可数序稠密集 $D=\{d_1,d_2,\ldots\}$，逐个插入有理数。插入 $d_n$ 时，若已有点在其下方和上方，取两者有理数坐标的中点；若只有一侧有界，则向该侧取足够远的有理数。每一步只改变新点坐标，不改变已有点之间的顺序。由归纳法，最终映射保持严格序。

\begin{proposition}[中点构造保持序]
若 $a<b$ 且 $a,b\in\Q$，则 $(a+b)/2\in\Q$ 且 $a<(a+b)/2<b$。
\end{proposition}
\begin{proof}
有理数对加法和除以非零整数封闭，故中点为有理数；两边不等式由 $b-a>0$ 直接得到。
\end{proof}

\begin{example}[三次插入]
先把 $d_1$ 映射为 $0$，把 $d_2\succ d_1$ 映射为 $1$。若第三个元素满足 $d_2\succ d_3\succ d_1$，则取 $g(d_3)=1/2$。第四个元素若位于 $d_3$ 与 $d_2$ 之间，则取 $3/4$。每次取中点，都保持已有坐标的相对顺序。
\end{example}

\subsection{可加表示的检验}
若偏好具有坐标独立性与矩形条件，可写为 $u(x_1,x_2)=v_1(x_1)+v_2(x_2)$。检验时先固定第二坐标比较第一坐标，再固定第一坐标比较第二坐标，最后用矩形条件排除交叉项。该步骤解释了为什么“每一维都可单独排序”仍不足以推出可加性：还需要不同坐标变化之间的替代关系一致。

\begin{remark}
表示定理的证明不仅回答存在性，也揭示限制来源：有限集依赖排序，可数集依赖稠密插入，连续消费空间依赖拓扑连续性，可加表示还依赖跨坐标一致性。不同结论不能互相替代。
\end{remark}
